% !PRINTED 310
\weberpage{310}

Wir wollen zunächst die Transformationsgleichung \S\,72, (8)'
in der jetzt in Übereinstimmung mit dem vorigen Paragraphen
$v$ für $x$ gesetzt ist:
\weq{1}{v^6 + 10\,v^3 - \gamma_2 v + 5 = 0}
anwenden. Hierin ist, wenn
\[
c \equiv 0 \ (\mathrm{mod}\ 12), \quad z \equiv c \ (\mathrm{mod}\ 5)
\]
ist,
\weq{2}{v_\infty = 5\left(\frac{\eta(5\omega)}{\eta\,\omega}\right)^2, \quad v_z = \left(\frac{\eta\left(\dfrac{\omega+c}{5}\right)}{\eta(\omega)}\right)^2.} % [sic] Die Quelle druckt im ersten Nenner \eta\omega ohne Funktionsklammern; der parallele zweite Nenner druckt \eta(\omega).
Nehmen wir dann
\weq{3}{\begin{aligned}
\sqrt{5}\,w_z &= (v_\infty - v_z)(v_{z+1} - v_{z-1})(v_{z-2} - v_{z+2}) \\
&= (v_\infty - v_c)(v_{c+12} - v_{c-12})(v_{c+24} - v_{c-24}), \\
z &= 0, 1, 2, 3, 4, \quad c = 0, \pm 12, \pm 24,
\end{aligned}}
so sind nach der Schlußbemerkung des vorigen Paragraphen die
$w_0$, $w_1$, $w_2$, $w_3$, $w_4$ die Wurzeln einer Gleichung 5ten Grades,
deren Koeffizienten rational aus $\gamma_2$ und rationalen Zahlen gebildet
 sind.

Beachtet man aber die Relationen [\S\,54, (14)]:
\[
\gamma_2(\omega+1) = e^{-\frac{2\pi i}{3}} \gamma_2(\omega), \quad \gamma_2\left(-\frac{1}{\omega}\right) = \gamma_2(\omega),
\]
so folgt aus der Form von (1), daß durch die Vertauschung
$(\omega, \omega+1)$ die Wurzeln $v$, abgesehen von einer Umstellung, den
Faktor $e^{\frac{2\pi i}{3}}$ annehmen, während sie durch $\left(\omega, -\dfrac{1}{\omega}\right)$ ungeändert
bleiben.

Die $w_z$ vertauschen sich also nur untereinander, und
ihre symmetrischen Funktionen bleiben \gesperrt{ungeändert} und hängen
daher rational nur von der Invariante $j(\omega)$ ab. Die Koeffizienten
in der Gleichung für $w$ können überdies für kein endliches $j(\omega)$
unendlich werden und sind demnach \gesperrt{ganze Funktionen}
von $j(\omega)$.

Ein weiterer Aufschluß über diese Koeffizienten ergiebt sich
durch die Entwickelung nach steigenden Potenzen von $q$. Die
Entwickelungen der $v$ haben nach (2) die Anfänge
\[
v_\infty = 5\,q^{\frac{2}{3}} + \cdots, \quad v_z = e^{\frac{c}{12}\frac{2\pi i}{5}} q^{-\frac{2}{15}} + \cdots,
\]
und daraus folgt mit Benutzung der bekannten Gleichung
\[
4\sin\frac{2\pi}{5}\cdot\sin\frac{4\pi}{5} = \sqrt{5}
\]
%CONT>>
